% Complete conditional-expectation proofs; kernels are a separate construction.

\begin{annexitem}{annex:p15.conditional-existence}{thm:p15-conditional-existence}{Integrability, uniqueness and the actual density construction}
\begin{proof}
Suppose \(Z\) satisfies the defining conditions, and take
\(A=\{Z>0\}\in\mathcal G\).
The event identity on \(A\) and on \(A^c\) gives
\[
          \int_A Z\,dP=\int_A X\,dP\leq\int_A|X|\,dP,\qquad
          \int_{A^c}(-Z)\,dP=\int_{A^c}(-X)\,dP
                                      \leq\int_{A^c}|X|\,dP.
\]
Both left sides are nonnegative, and the right sides are finite.
Adding proves \(E|Z|\leq E|X|<\infty\).
If \(Z'\) is another version, for
\(A_n=\{Z\geq Z'+1/n\}\in\mathcal G\), signed linearity gives
\[
                       0=\int_{A_n}(Z-Z')\,dP\geq P(A_n)/n.
\]
Thus each \(A_n\) is null, and their union \(\{Z>Z'\}\) is null.
Interchanging \(Z,Z'\) proves equality almost everywhere.

For existence, restrict \(P\) to \(\mathcal G\). On that measurable
space the two set functions
\[
                       \nu_\pm(A)=\int_A X^\pm\,dP
\]
are finite measures, by the nonnegative integral construction and
integrability of \(X\).
Each is absolutely continuous with respect to \(P|_{\mathcal G}\).
Apply the proved Radon--Nikodym theorem on this exact space.
The densities \(Z_+,Z_-\) are \(\mathcal G\)-measurable, nonnegative
and integrable, with integrals \(EX^+,EX^-\).
Their infinite-value sets are measurable in \(\mathcal G\) and null.
Set both to zero on the union of those sets, obtaining real
\(\mathcal G\)-measurable representatives.
Then \(Z=Z_+-Z_-\) is integrable and has, for every \(A\in\mathcal G\),
\[
                       \int_AZ\,dP=\nu_+(A)-\nu_-(A)=\int_AX\,dP.
\]
No completion of the conditioning sigma-algebra was used.
\end{proof}
\end{annexitem}

\begin{annexitem}{annex:p15.conditional-basic}{thm:p15-conditional-basic}{Testing each claim on the conditioning events}
\begin{proof}
The linear combination \(aC_{\mathcal G}X+bC_{\mathcal G}Y\)
is measurable and integrable. Its integral on any conditioning event
is \(a\int_AX+b\int_AY=\int_A(aX+bY)\).
Uniqueness proves linearity.
Choosing \(A=\Omega\) in the definition gives the expectation identity.

If \(X\leq Y\) almost everywhere, put
\(A_n=\{C_{\mathcal G}X\geq C_{\mathcal G}Y+1/n\}\in\mathcal G\).
Then
\[
       P(A_n)/n\leq\int_{A_n}(C_{\mathcal G}X-C_{\mathcal G}Y)\,dP
                            =\int_{A_n}(X-Y)\,dP\leq0.
\]
Their union is null, proving order. Apply it to
\(-|X|\leq X\leq|X|\) and use linearity to obtain the absolute-value bound.

If \(X\) is already \(\mathcal G\)-measurable, it itself satisfies
both defining conditions. If \(X\) is independent of \(\mathcal G\),
then for every \(A\in\mathcal G\) the variables \(X,\mathbf1_A\)
are independent. The proved expectation factorization gives
\(\int_AX\,dP=EX\,P(A)\), exactly the integral of the constant \(EX\)
on that event. The constant is measurable and integrable, proving the
independence assertion.

For locality, put
\[
 A_n=B\cap\{C_{\mathcal G}X\geq C_{\mathcal G}Y+1/n\}.
\]
These are conditioning events contained in \(B\), where \(X-Y=0\)
almost everywhere. Their integral identity gives \(P(A_n)=0\).
The reversed inequality is treated by interchanging \(X,Y\).
The conclusion holds on \(B\); neither calculation asserts that an
event outside \(B\) is null.
Finally condition \(0\leq\mathbf1_A\leq1\), and use the fact that
the constants 0 and 1 are their own conditional expectations.
\end{proof}
\end{annexitem}

\begin{annexitem}{annex:p15.conditional-limits}{thm:p15-conditional-limits}{Nonnegative versions and the conditional limit arguments}
\begin{proof}
\paragraph{Extension by truncation.}
For \(X\geq0\), choose versions \(Z_n=C_{\mathcal G}(X\wedge n)\).
Order and positivity show \(0\leq Z_n\leq Z_{n+1}\) almost everywhere.
The union of their exceptional sets is a \(\mathcal G\)-measurable
null set, since it is described by these \(\mathcal G\)-measurable
functions. Set every \(Z_n\) to zero there.
Their pointwise increasing limit \(Z\in[0,\infty]\) is
\(\mathcal G\)-measurable. Monotone convergence on each event gives
\[
                   \int_AZ\,dP
                     =\lim_n\int_AZ_n\,dP
                     =\lim_n\int_A(X\wedge n)\,dP
                     =\int_AX\,dP.
\]

For uniqueness among nonnegative extended versions \(Z,W\), consider
\[
                  A_{mn}=\{Z\leq m,\ W\geq Z+1/n\}\in\mathcal G.
\]
The integral of \(Z\) there is at most \(m\), and equality of the
two defining integrals makes the integral of \(W\) finite there too.
Their difference is therefore legitimate and at least \(P(A_{mn})/n\).
It must be zero, so \(A_{mn}\) is null.
Their union is the event \(\{Z<W,\ Z<\infty\}\).
Interchanging \(Z,W\) covers the other possible inequality; both
infinite values already agree. This proves uniqueness and compatibility
with the nonnegative integrable definition.

This same bounded-level comparison proves order for the nonnegative
extension. If \(X\leq Y\), for integers \(m,n\geq1\) take
\[
 A_{mn}=\{C_{\mathcal G}Y\leq m,\
             C_{\mathcal G}X\geq C_{\mathcal G}Y+1/n\}.
\]
The event integral for \(Y\) is at most \(m\), and the event identity
and \(X\leq Y\) bound the integral for \(X\) by the same finite number.
Their difference is therefore defined, is at most zero, and is at
least \(P(A_{mn})/n\). Hence \(P(A_{mn})=0\).
Their countable union covers every possible strict reversed inequality
where \(C_{\mathcal G}Y\) is finite; a strict reversal is impossible
where it is infinite. This proves order on the whole space almost everywhere.

\paragraph{Monotone convergence and Fatou.}
For \(X_n\uparrow X\), order allows the versions \(C_{\mathcal G}X_n\)
to be increasing outside one conditioning-null set.
Their limit has event integrals \(\lim_n\int_AX_n=\int_AX\),
so uniqueness identifies it with \(C_{\mathcal G}X\).
Almost-everywhere hypotheses give the same identities because changing
the integrands on a measurable \(P\)-null set changes no integral.

For arbitrary nonnegative \(X_n\), set \(Y_n=\inf_{k\geq n}X_k\).
Then \(Y_n\uparrow\liminf_kX_k\).
The countably many order comparisons give, outside one null set,
\[
                       C_{\mathcal G}Y_n
                               \leq\inf_{k\geq n}C_{\mathcal G}X_k.
\]
Let \(n\to\infty\) and use conditional monotone convergence to obtain
the stated Fatou inequality.

\paragraph{Dominated convergence.}
If \(|X_n|\leq Y\) with \(EY<\infty\) and \(X_n\to X\) almost
everywhere, then \(|X|\leq Y\) almost everywhere and \(X\) is integrable.
Conditional Fatou applied to \(Y+X_n\) gives
\[
                   C_{\mathcal G}Y+C_{\mathcal G}X
                      \leq C_{\mathcal G}Y+\liminf_nC_{\mathcal G}X_n.
\]
These quantities are finite where subtraction is used, because
all conditional absolute values are bounded by \(C_{\mathcal G}Y\),
which is finite almost everywhere. Applying Fatou to \(Y-X_n\)
likewise gives \(\limsup_nC_{\mathcal G}X_n\leq C_{\mathcal G}X\).
This proves almost-everywhere convergence.
Linearity and the integrability bound already proved give
\[
            E|C_{\mathcal G}X_n-C_{\mathcal G}X|\leq E|X_n-X|\to0,
\]
where the last step is ordinary dominated convergence.
Every adjustment of conditional versions made here occurred on a
measurable null set in \(\mathcal G\); this preserves the required
measurability without completing that sigma-algebra.
That a change on an \(\mathcal F_0\)-null set outside \(\mathcal G\)
can destroy this measurability is proved, with a witness, in
\cref{annex:p15-w02b-p15-008-null-modification}; the decreasing variant
is derived from item~1 of this theorem in
\cref{annex:p15-w02b-p15-008-decreasing}.
\end{proof}
\end{annexitem}

\begin{annexitem}{annex:p15.conditional-jensen}{thm:p15-conditional-jensen}{Rational affine minorants and a countable exceptional union}
\begin{proof}
\paragraph{The countable approximation.}
The supporting-line proof in \cref{thm:p15-jensen-integral} supplies,
at every real \(x\), an affine function \(L_x(t)=a_xt+b_x\) that
is everywhere at most \(\varphi(t)\) and equals it at \(x\).
If all these selected slopes were one number \(a\), then the
supporting inequalities from \(x\) to \(y\) and back would give
\(\varphi(y)-\varphi(x)=a(y-x)\) for every \(x,y\).
Thus a nonaffine \(\varphi\) has two such lines of different slopes.

Fix \(x\) and \(\varepsilon>0\), and choose a supporting line \(L_x\).
Of the two lines just found, at least one, say \(L'\), has a
different slope. For \(0<\theta<1\), the convex combination
\[
                             L_\theta=(1-\theta)L_x+\theta L'
\]
is still a global affine minorant.
Choose \(\theta_0>0\) sufficiently small that
\(L_\theta(x)>\varphi(x)-\varepsilon/2\) for \(0<\theta<\theta_0\).
Its slopes range through a nonempty real interval as \(\theta\)
ranges through this interval. Density of the rationals gives a
\(\theta\) for which the slope is rational.
Lower its real intercept to a rational number by less than
\(\varepsilon/2\). The resulting rational-coefficient line is still
a global minorant and exceeds \(\varphi(x)-\varepsilon\) at \(x\).
Since every minorant is at most \(\varphi\), this proves the exact
supremum representation by the countable set of rational pairs.

For an affine \(\varphi(t)=at+b\), the conditional inequality is
an equality by linearity. If \(a\) is irrational, a line with rational
slope has a different slope and eventually exceeds \(\varphi\) in
one direction, regardless of its intercept. This proves the stated
reason that the rational-minorant construction needs a separate affine case.

\paragraph{Taking the conditional supremum.}
In the nonaffine case, for each rational minorant \(at+b\leq\varphi(t)\),
conditional order and linearity give
\[
                  aC_{\mathcal G}X+b
                       \leq C_{\mathcal G}(\varphi(X))
                             \quad\text{almost everywhere}.
\]
Both variables whose conditional expectations occur are integrable
by the hypotheses. There are only countably many pairs \((a,b)\);
outside their union of exceptional null sets all inequalities hold
at once. Taking the pointwise supremum there gives
\(\varphi(C_{\mathcal G}X)\leq C_{\mathcal G}(\varphi(X))\).
This proves the source's rational-coefficient argument in full.
\end{proof}
\end{annexitem}

\begin{annexitem}{annex:p15.conditional-lp}{thm:p15-conditional-lp}{The norm properties and the contraction inequalities}
\begin{proof}
Null-set replacement does not change any displayed quantity.
For finite \(p\), positivity, scalar homogeneity and the fact that
zero norm means the zero almost-everywhere class follow from the
nonnegative integral laws. The triangle inequality at \(p=1\)
follows from \(|X+Y|\leq|X|+|Y|\).
For \(p>1\), the real function \(|t|^p\) is convex: its derivative
\(p\,\operatorname{sgn}(t)|t|^{p-1}\), where
\(\operatorname{sgn}(t)\) is \(-1,0,1\) according as \(t<0,t=0,t>0\),
is nondecreasing; integrating that monotonicity gives the secant
inequalities for convexity. In particular
\[
                    |X+Y|^p\leq2^{p-1}(|X|^p+|Y|^p),
\]
so \(X+Y\in L^p\).
Let \(q=p/(p-1)\). Applying the proved H\"older inequality twice gives
\[
 \begin{aligned}
 \|X+Y\|_p^p
   &\leq E\bigl((|X|+|Y|)|X+Y|^{p-1}\bigr)\\
   &\leq(\|X\|_p+\|Y\|_p)\,\bigl\||X+Y|^{p-1}\bigr\|_q\\
   &=(\|X\|_p+\|Y\|_p)\|X+Y\|_p^{p-1}.
 \end{aligned}
\]
If the last norm is zero the triangle inequality is immediate;
otherwise division gives it.

For an essentially bounded variable, the infimum defining
\(\|X\|_\infty\) is itself an almost-everywhere bound:
choose bounds approaching that infimum and remove the countable
union of their exceptional sets. The limit is then a bound too.
Positivity, homogeneity, definiteness and the triangle inequality
follow from the corresponding pointwise properties outside such
null sets. This completes all the norm assertions.

If \(1\leq p<\infty\) and \(X\in L^p\), H\"older against the constant
one shows \(X\in L^1\). Conditional Jensen, followed by the defining
expectation identity, gives
\[
                    E|C_{\mathcal G}X|^p
                       \leq E(C_{\mathcal G}|X|^p)=E|X|^p.
\]
For \(p=\infty\), condition the bound \(|X|\leq\|X\|_\infty\)
and use \(|C_{\mathcal G}X|\leq C_{\mathcal G}|X|\).
This gives the corresponding essential bound.
Linearity applied to \(X-Y\) proves each difference inequality,
and hence preservation of the stated norm convergence.
\end{proof}
\end{annexitem}

\begin{annexitem}{annex:p15.conditional-tower}{thm:p15-conditional-tower}{Nested conditioning checked in both directions}
\begin{proof}
The \(\mathcal G\)-measurable variable \(C_{\mathcal G}X\) is
also \(\mathcal H\)-measurable and integrable, so conditioning it
on \(\mathcal H\) leaves it unchanged. This proves the first identity.
For the second, \(C_{\mathcal G}X\) is measurable and integrable
for the outer conditioning space, and for any \(A\in\mathcal G\)
the two defining identities give
\[
                 \int_A C_{\mathcal G}X\,dP
                    =\int_AX\,dP
                    =\int_A C_{\mathcal H}X\,dP.
\]
It is therefore a version of \(C_{\mathcal G}(C_{\mathcal H}X)\).
If \(C_{\mathcal H}X\) itself has a \(\mathcal G\)-measurable
version, its event identity on all \(\mathcal H\) events holds
in particular on \(\mathcal G\), proving the final assertion.
\end{proof}
\end{annexitem}

\begin{annexitem}{annex:p15.conditional-pullout}{thm:p15-conditional-pullout}{The full nonnegative approximation and signed decomposition}
\begin{proof}
First take a nonnegative extended \(X\) and \(Z=\mathbf1_B\)
for \(B\in\mathcal G\). For every \(A\in\mathcal G\),
\[
                 \int_A\mathbf1_B C_{\mathcal G}X\,dP
                     =\int_{A\cap B}X\,dP
                     =\int_A\mathbf1_BX\,dP.
\]
Nonnegative uniqueness proves the identity.
Finite nonnegative linear combinations give the identity for
nonnegative simple \(Z\): event integrals add in the nonnegative
extended sense, so the sum is the required version.
For arbitrary nonnegative extended \(\mathcal G\)-measurable \(Z\),
choose nonnegative simple \(Z_n\uparrow Z\).
Monotone convergence on both sides of the event identity yields
\[
                       \int_A ZC_{\mathcal G}X\,dP=\int_AZX\,dP.
\]
The product is measurable and nonnegative, with zero times infinity
defined as zero. Nonnegative uniqueness proves the complete extended rule.

Now take real \(Z\), integrable \(X\), and \(E|ZX|<\infty\).
The nonnegative rule applied to \(|Z|,|X|\) gives
\[
                   E\bigl(|Z|C_{\mathcal G}|X|\bigr)=E|ZX|<\infty.
\]
Since \(|ZC_{\mathcal G}X|\leq |Z|C_{\mathcal G}|X|\), the right
side of the claimed signed identity is integrable.
Expand \(ZX=(Z^+-Z^-)(X^+-X^-)\).
Each of the four nonnegative products is at most \(|ZX|\),
and its conditional counterpart \(Z^\pm C_{\mathcal G}X^\pm\)
has the same finite integral by the nonnegative rule.
Applying signed linearity to these four integrable terms gives
\[
                 C_{\mathcal G}(ZX)
                      =(Z^+-Z^-)(C_{\mathcal G}X^+-C_{\mathcal G}X^-)
                      =ZC_{\mathcal G}X .
\]
This also verifies the asserted right-side integrability without
assuming \(E|Z|<\infty\).
\end{proof}
\end{annexitem}

\begin{annexitem}{annex:p15.conditional-partitions}{thm:p15-conditional-partitions}{All cells, their null cases and arbitrary unions}
\begin{proof}
Every member of \(\mathcal G=\sigma(A_i:i)\) is a union of cells.
Indeed, the collection of all unions of cells is a sigma-algebra
containing the cells, and each such union is a finite or countable
union and hence belongs to \(\mathcal G\).
Define \(Z\) by the displayed constant on each positive cell and
zero on each null cell. It is \(\mathcal G\)-measurable, since
every inverse image of a Borel set is a union of cells.
Moreover,
\[
             E|Z|
                =\sum_{i:P(A_i)>0}\left|\int_{A_i}X\,dP\right|
                \leq\sum_i\int_{A_i}|X|\,dP=E|X|.
\]
Thus it is integrable. On each positive cell its integral equals
\(\int_{A_i}X\) by multiplication by that cell's probability.
On a null cell both integrals vanish.
Absolute convergence of the two signed integral sums then proves
the event identity on every union of cells, hence every event in
\(\mathcal G\). Uniqueness gives the formula.

The union of the null cells is measurable in \(\mathcal G\)
and null, because there are only countably many cells.
Other finite constants on those cells preserve measurability and
all integrals. The partition with the single cell \(\Omega\)
gives the constant \(EX\), and restricting to a finite partition
recovers the earlier finite formula.
\end{proof}
\end{annexitem}

\begin{annexitem}{annex:p15.conditional-factorization}{thm:p15-conditional-factorization}{Measurable lifts without assumptions on the observation image}
\begin{proof}
The collection \(\{Y^{-1}(B):B\in\mathcal T\}\) is a sigma-algebra
on \(\Omega\): preimages commute with complements and countable unions,
and the preimage of \(T\) is \(\Omega\).
It is therefore exactly \(\sigma(Y)\).

For a nonnegative simple \(Z=\sum_{i=1}^m a_i\mathbf1_{A_i}\)
with disjoint \(\sigma(Y)\) cells covering \(\Omega\), choose
\(B_i\in\mathcal T\) with \(Y^{-1}(B_i)=A_i\).
Disjointify the \(B_i\) in their listed order to obtain \(D_i\).
Their preimages remain \(A_i\), since the original preimages were
disjoint. Define \(h=\sum_i a_i\mathbf1_{D_i}\), with value zero
on the remaining part of \(T\). Then \(h\) is measurable and
\(h(Y)=Z\) everywhere.

For a nonnegative extended \(Z\), choose its increasing nonnegative
simple approximations \(Z_n\), and lift each to a nonnegative
measurable \(h_n\) as above. The lifts need not be increasing
away from the observation image. Set \(h=\limsup_nh_n\), which
is measurable and nonnegative extended on all of \(T\).
At every actual observation, \(h(Y)=\lim_nZ_n=Z\).
For a real \(Z\), lift \(Z^+,Z^-\) separately.
On the measurable set where both lifted functions are finite,
take their difference; assign zero elsewhere. At every point
of \(Y(\Omega)\) both are finite, so this measurable real function
still composes with \(Y\) to give exactly \(Z\).
This construction neither tests nor assumes measurability of
the image \(Y(\Omega)\).

Apply the result to a selected \(\sigma(Y)\)-measurable version
of \(E(X\mid Y)\). If \(h(Y)\) and \(k(Y)\) give the same version
almost everywhere, their inequality event has probability
\[
             P(h(Y)\ne k(Y))=\mu_Y(\{t:h(t)\ne k(t)\}),
\]
where \(\mu_Y\) is the law of \(Y\).
This proves the precise almost-everywhere uniqueness on the observation
space and the qualification on assigning values to null observations.
\end{proof}
\end{annexitem}

\begin{annexitem}{annex:p15.conditional-density-formulas}{thm:p15-conditional-density-formulas}{Finite signed sections and their event identities}
\begin{proof}
\paragraph{Joint real density.}
Tonelli makes \(m(y)=\int f(x,y)\,dx\) measurable and gives
\(\int m(y)\,dy=1\). It is the marginal density of \(Y\), by
integrating over a rectangle and then over each event \(\{Y\in B\}\).
Likewise \(a(y)=\int|g(x)|f(x,y)\,dx\) is measurable and
\(\int a(y)\,dy=E|g(X)|<\infty\).
Thus the sets \(\{m=\infty\}\) and \(\{a=\infty\}\) are Lebesgue-null,
and hence null under the \(Y\)-law, which has density \(m\).
The set \(\{m=0\}\) also has \(Y\)-mass zero.
Consequently \(P(Y\in D)=1\).

On \(D\), the signed numerator is the difference of two finite
measurable nonnegative section integrals. Dividing by the positive
finite denominator there and extending by zero gives a measurable
real \(h\). Moreover
\[
                    E|h(Y)|=\int_D|h(y)|m(y)\,dy
                                   \leq\int_Da(y)\,dy<\infty.
\]
For any Borel \(B\), Fubini and the null exceptional sections give
\[
 \begin{aligned}
 E(h(Y)\mathbf1_{\{Y\in B\}})
   &=\int_{B\cap D}\int g(x)f(x,y)\,dx\,dy\\
   &=E(g(X)\mathbf1_{\{Y\in B\}}).
 \end{aligned}
\]
Every event in \(\sigma(Y)\) is such a preimage, by the preceding
factorization proof's initial sigma-algebra identity.
This proves both defining conditions and hence the formula.

\paragraph{An independent coordinate.}
Independence gives joint law \(\mu\times\nu\) by the earlier
product-law theorem. Fubini applied to the integrable \(\phi\)
gives a measurable finite section integral on a full-\(\mu\)-measure
set and the stated zero extension elsewhere.
It also gives \(\int|h|\,d\mu\leq\int|\phi|\,d(\mu\times\nu)\).
For any \(A=\{X\in B\}\in\sigma(X)\),
\[
 \begin{aligned}
 E(\phi(X,Y)\mathbf1_A)
   &=\int_{B\times T}\phi(s,t)\,d(\mu\times\nu)(s,t)\\
   &=\int_Bh(s)\,d\mu(s)=E(h(X)\mathbf1_A).
 \end{aligned}
\]
The function \(h(X)\) is measurable from \(X\) and integrable.
Uniqueness proves the second assertion on precisely the stated
product measurable domain.
\end{proof}
\end{annexitem}

\begin{annexitem}{annex:p15.conditional-l2-space}{thm:p15-conditional-l2-space}{A complete norm limit with a measurable representative}
\begin{proof}
Cauchy--Schwarz makes \(E(XY)\) finite for two square-integrable
variables. Integral linearity and symmetry make it a symmetric
bilinear form, and \(E(X^2)=0\) holds exactly for the zero
almost-everywhere class. It is therefore an inner product, inducing
the already proved norm \(\|X\|_2\).

To prove completeness for either sigma-algebra \(\mathcal A\), let
\((X_n)\) be Cauchy in that norm, with real \(\mathcal A\)-measurable
representatives. Choose increasing indices \(n_j\) so that
\[
                  \|X_{n_{j+1}}-X_{n_j}\|_2\leq2^{-j}.
\]
Such a subsequence follows from the Cauchy property by choosing each
index beyond the tail bound needed for that stage.
Put \(D_j=X_{n_{j+1}}-X_{n_j}\).
For finite sums, the norm triangle inequality gives
\[
                  \left\|\sum_{j=1}^N|D_j|\right\|_2
                                          \leq\sum_{j=1}^N2^{-j}\leq1.
\]
Monotone convergence of their squares proves that
\(H=\sum_{j\geq1}|D_j|\) has \(EH^2\leq1\).
Its infinite-value set is therefore \(\mathcal A\)-measurable and null.
On its complement, the series \(\sum_jD_j\) converges absolutely.
Define \(X=X_{n_1}+\sum_jD_j\) there, and set \(X=0\) on that null set.
It is a real \(\mathcal A\)-measurable variable, with
\(|X|\leq|X_{n_1}|+H\) almost everywhere, hence is in \(L^2\).

The same finite-sum and monotone-limit argument for each tail gives
\[
                 \left\|\sum_{j\geq m}|D_j|\right\|_2
                                      \leq\sum_{j\geq m}2^{-j}.
\]
Since \(|X-X_{n_m}|\) is bounded by this tail almost everywhere,
\(X_{n_m}\to X\) in norm. The original Cauchy property and the norm
triangle inequality then give \(X_n\to X\) for the whole sequence.
This proves completeness without a pointwise limit on a nonmeasurable
exceptional subset.

Apply the argument first to \(\mathcal F_0\) and then to \(\mathcal G\).
The \(\mathcal G\)-classes embed linearly and isometrically into the
\(\mathcal F_0\)-classes. If a sequence from that subspace converges
in the ambient norm, it is Cauchy in the identical subspace norm,
and the proved completeness gives a \(\mathcal G\)-measurable limit
there. Uniqueness of norm limits identifies that class with the
ambient limit. Thus the subspace is closed.
\end{proof}
\end{annexitem}

\begin{annexitem}{annex:p15.conditional-l2-projection}{thm:p15-conditional-l2-projection}{Orthogonality and the full Pythagorean identity}
\begin{proof}
Contraction gives \(M=C_{\mathcal G}X\in L^2(\mathcal G)\).
For \(Z\in L^2(\mathcal G)\), both \(ZX\) and \(ZM\) are
integrable by Cauchy--Schwarz.
The measurable-multiplier rule gives \(C_{\mathcal G}(ZX)=ZM\).
Taking expectations yields \(E(ZX)=E(ZM)\), and hence
\(E[Z(X-M)]=0\).

For \(Y\in L^2(\mathcal G)\), choose \(Z=M-Y\) in this identity.
Expanding a square and using the zero cross term gives
\[
 E(X-Y)^2
   =E\bigl((X-M)+(M-Y)\bigr)^2
   =E(X-M)^2+E(M-Y)^2 .
\]
All three squares and the cross term are integrable by the stated
second moments and Cauchy--Schwarz.
The last summand is nonnegative and vanishes exactly when \(Y=M\)
almost everywhere. This proves minimality and its precise uniqueness.
If any real \(\mathcal G\)-measurable \(Y\) has finite error, then
\[
                          Y^2\leq2(X-Y)^2+2X^2
\]
makes it square integrable. A competitor with infinite error cannot
improve the finite error of \(M\). This proves the assertion for
all real measurable competitors.
\end{proof}
\end{annexitem}

\begin{annexitem}{annex:p15.conditional-variance}{thm:p15-conditional-variance}{Residual variance and nested prediction errors}
\begin{proof}
Put \(M=C_{\mathcal G}X\).
The products \(XM\) and \(M^2\) are integrable, so conditional
linearity and the multiplier rule give
\[
 \begin{aligned}
 C_{\mathcal G}((X-M)^2)
   &=C_{\mathcal G}(X^2)-2M C_{\mathcal G}X+M^2\\
   &=C_{\mathcal G}(X^2)-M^2 .
 \end{aligned}
\]
The left side is nonnegative by conditional order. Its expectation
is \(EX^2-EM^2\).
Since \(EM=EX\), adding \(\operatorname{Var}(M)=EM^2-(EX)^2\)
proves total variance.

If \(\mathcal G\subseteq\mathcal H\), then
\(C_{\mathcal G}X\in L^2(\mathcal H)\).
Apply the \(\mathcal H\)-projection identity with that variable
as competitor to obtain exactly
\[
 E(X-C_{\mathcal G}X)^2
    =E(X-C_{\mathcal H}X)^2
             +E(C_{\mathcal H}X-C_{\mathcal G}X)^2 .
\]
Dropping the last nonnegative term proves the information comparison.
\end{proof}
\end{annexitem}

\begin{annexitem}{annex:p15.conditional-inequalities}{thm:p15-conditional-inequalities}{Indicator bounds, rational quadratics and Bayes}
\begin{proof}
The pointwise inequality
\(\mathbf1_{\{|X|\geq a\}}\leq X^2/a^2\) and conditional
order give conditional Chebyshev immediately.

For conditional Cauchy--Schwarz, set
\[
        A=C_{\mathcal G}(X^2),\quad
        B=C_{\mathcal G}(XY),\quad
        D=C_{\mathcal G}(Y^2).
\]
They are finite almost everywhere: the squares are integrable,
and \(XY\) is integrable by ordinary Cauchy--Schwarz.
For every rational \(t\), conditional positivity and linearity give
\[
                     A+2tB+t^2D=C_{\mathcal G}((X+tY)^2)\geq0
                                      \quad\text{almost everywhere}.
\]
Remove the single null set combining the countably many rational
exceptions and the three finiteness exceptions.
At every remaining point, this real polynomial is nonnegative for
all rational \(t\), hence for every real \(t\) by continuity.
If \(D>0\), evaluating it at \(t=-B/D\) gives \(B^2\leq AD\).
If \(D=0\), a nonconstant affine polynomial could not be nonnegative
for every real \(t\), so \(B=0\), and the same inequality follows.
This supplies all cases of the conditional bound.

For Bayes, let \(h=P(A\mid\mathcal G)\).
Its defining identity gives
\(\int_Bh\,dP=P(A\cap B)\) for \(B\in\mathcal G\), and
\(\int_\Omega h\,dP=P(A)>0\).
Their ratio is exactly \(P(B\mid A)\).
For a countable partition, \(h\) has value \(P(A\mid A_i)\)
on each positive cell. Substitution gives numerator
\(P(A\mid A_i)P(A_i)\) for \(B=A_i\) and denominator
\(\sum_jP(A\mid A_j)P(A_j)\), with null cells contributing zero.
\end{proof}
\end{annexitem}

\begin{annexitem}{annex:p15.conditional-random-sums}{thm:p15-conditional-random-sums}{Integrability before conditioning on the random count}
\begin{proof}
For each fixed \(n\geq1\),
\[
                    \left(\sum_{i=1}^nY_i\right)^2
                                \leq n\sum_{i=1}^nY_i^2
\]
by finite Cauchy--Schwarz. At \(n=0\), both sides are zero.
The variables \(N\) and every finite vector of \(Y_i\)'s are independent.
Decomposing into its countably many \(N\)-events and using nonnegative
sum interchange therefore gives
\[
 ES^2
    =\sum_{n\geq0}E\left[\mathbf1_{\{N=n\}}
                        \left(\sum_{i=1}^nY_i\right)^2\right]
    \leq\sum_{n\geq0}P(N=n)n^2EY_1^2
    =EN^2\,EY_1^2<\infty.
\]
Thus the conditional moment and variance formulas are applicable.
Also \(N\) and \(S\) are integrable by their finite second moments.

For a fixed \(n\), the finite independent-sum formulas give
mean \(n\mu\) and second moment \(n\sigma^2+n^2\mu^2\).
Independence from \(\{N=n\}\) shows these are the corresponding
conditional values on every positive-probability count cell.
The countable-partition formula consequently gives versions
\[
              E(S\mid N)=\mu N,\qquad
              E(S^2\mid N)=\sigma^2N+\mu^2N^2.
\]
The displayed values may also be used on null count cells, preserving
measurability and all integrals.
Thus \(\operatorname{Var}(S\mid N)=\sigma^2N\).
The expectation identity gives \(ES=\mu EN\), and total variance
gives
\(\operatorname{Var}(S)=\sigma^2EN+\operatorname{Var}(\mu N)\),
which is exactly the stated expression, including the empty-sum case.
\end{proof}
\end{annexitem}

\begin{annexitem}{annex:p15.conditional-same-law}{thm:p15-conditional-same-law}{The shifted-sign proof for integrable variables}
\begin{proof}
Let \(M=C_{\mathcal G}X\) and assume \(M,X\) have the same law.
For any fixed real \(c\), both \(X-c\) and \(M-c\) are integrable
and have equal absolute first moments.
Conditional Jensen for absolute value gives
\[
                         C_{\mathcal G}|X-c|\geq|M-c|.
\]
The expectation of their nonnegative difference is zero by equality
of laws. Hence they agree almost everywhere.

On \(A_+=\{M>c\}\in\mathcal G\), the defining integral identity yields
\[
 E\bigl((|X-c|-(X-c))\mathbf1_{A_+}\bigr)
   =E\bigl((C_{\mathcal G}|X-c|-(M-c))\mathbf1_{A_+}\bigr)=0.
\]
The integrand on the left is nonnegative, so \(X\geq c\)
almost everywhere on \(A_+\).
The corresponding calculation with
\(|X-c|+(X-c)\) on \(A_-=\{M<c\}\) gives \(X\leq c\) there.
On \(A_0=\{M=c\}\), the equality
\(C_{\mathcal G}|X-c|=|M-c|=0\) and the event identity give
\(E(|X-c|\mathbf1_{A_0})=0\), hence \(X=c\) there.

These facts imply \(\{X>c\}\subseteq\{M>c\}\) up to a null set.
The two events have equal probabilities because \(X,M\) have the same law.
Their difference is therefore null, so their indicators agree
almost everywhere.
Apply this to every rational \(c\), and remove the countable union
of the exceptional sets. Distinct real values of \(X,M\) at any
remaining point would have a rational number strictly between them,
contradicting equality of the corresponding threshold indicators.
Thus \(X=M\) almost everywhere.

For the square-integrable assertion, orthogonality gives
\(E(X-M)^2=EX^2-EM^2\). Equality of the two moments makes this
zero, implying \(X=M\) almost everywhere without the same-law premise.
\end{proof}
\end{annexitem}

\begin{annexitem}{annex:p15.conditional-examples}{thm:p15-conditional-examples}{All finite tables and a freely assignable null value}
\begin{proof}
For the three-point uniform space, cellwise averaging gives the table
\[
 \begin{array}{c|ccccc}
 \omega&X&C_{\mathcal G_1}X&C_{\mathcal G_2}X&
 C_{\mathcal G_2}C_{\mathcal G_1}X&C_{\mathcal G_1}C_{\mathcal G_2}X\\\hline
 a&1&1/2&1&1/2&1/2\\
 b&0&1/2&0&1/4&1/2\\
 c&0&0&0&1/4&0
 \end{array}
\]
For \(C_{\mathcal G_2}C_{\mathcal G_1}X\), the second conditioning
averages \(1/2\) and 0 over \(\{b,c\}\), giving \(1/4\) on both outcomes.
For \(C_{\mathcal G_1}C_{\mathcal G_2}X\), the first conditioning
already leaves \(X\) unchanged, and the other partition then averages
its values at \(a,b\).
The countable-partition theorem verifies the event identities, since
every conditioning event is a union of those displayed cells.
The unequal iterated-expectation columns prove noncommutation without nesting.

For the symmetric two-point variable, \(EX=0\) and \(X^2=1\)
everywhere. Conditioning on the trivial sigma-algebra therefore
gives the two claimed constants, demonstrating failure of
multiplicativity.

For the uniform square, independence of the coordinates gives
\(E(X\mid Y)=EX=1/2\). More explicitly, for any Borel \(B\subseteq[0,1]\),
\[
                  E(X\mathbf1_{\{Y\in B\}})
                         =\frac12P(Y\in B),
\]
which is the defining integral identity.
The event \(\{Y=0\}\) is in \(\sigma(Y)\) and has probability zero.
Replacing the value of the Borel function \(h\) only at zero by
any finite real \(c\) preserves Borel measurability and changes
\(h(Y)\) only on this conditioning-null event.
Every defining integral is unchanged. This proves the version
claim directly, without dividing by \(P(Y=0)=0\).
\end{proof}
\end{annexitem}
